\QHeading{第1题　口算。}
\par\noindent


\begin{oralgrid}
\OralItem{1}{5.1+5.02}{10.12}
\OralItem{2}{2.1-2.01}{0.09}
\OralItem{3}{1.23+1.32}{2.55}
\OralItem{4}{7.31-7.01}{0.3}
\OralItem{5}{0.99-0.9}{0.09}
\OralItem{6}{4.2+2.99}{7.19}
\OralItem{7}{0.7-0.07}{0.63}
\OralItem{8}{6.82-4.52}{2.3}
\OralItem{9}{0.5+0.51}{1.01}
\OralItem{10}{2.7+5.2}{7.9}
\OralItem{11}{0.43-0.24}{0.19}
\OralItem{12}{4.02-2.7}{1.32}\end{oralgrid}
\par

\QHeading{口算思路补充（对应上表题号）}

\Answer{上面列出得数，下面按小题给出不同计算思路。计数单位法适用于小数，整数题可比较凑整、拆分或商不变等方法。}
\MethodTitle{第1小题：$5.1+5.02$}

\MethodBlock{口算思路一：按相同计数单位计算}{先看成若干个百分之一，按整数加减，最后还原成小数。}{\begin{aligned}[t]&\phantom{={}}5.1+5.02\\
&=(510+502)\div 100\\
&=1012\div 100\\
&=10.12\end{aligned}}
\MethodTitle{第2小题：$2.1-2.01$}

\MethodBlock{口算思路一：按相同计数单位计算}{先看成若干个百分之一，按整数加减，最后还原成小数。}{\begin{aligned}[t]&\phantom{={}}2.1-2.01\\
&=(210-201)\div 100\\
&=9\div 100\\
&=0.09\end{aligned}}
\MethodTitle{第3小题：$1.23+1.32$}

\MethodBlock{口算思路一：按相同计数单位计算}{先看成若干个百分之一，按整数加减，最后还原成小数。}{\begin{aligned}[t]&\phantom{={}}1.23+1.32\\
&=(123+132)\div 100\\
&=255\div 100\\
&=2.55\end{aligned}}
\MethodTitle{第4小题：$7.31-7.01$}

\MethodBlock{口算思路一：按相同计数单位计算}{先看成若干个百分之一，按整数加减，最后还原成小数。}{\begin{aligned}[t]&\phantom{={}}7.31-7.01\\
&=(731-701)\div 100\\
&=30\div 100\\
&=0.3\end{aligned}}
\MethodTitle{第5小题：$0.99-0.9$}

\MethodBlock{口算思路一：按相同计数单位计算}{先看成若干个百分之一，按整数加减，最后还原成小数。}{\begin{aligned}[t]&\phantom{={}}0.99-0.9\\
&=(99-90)\div 100\\
&=9\div 100\\
&=0.09\end{aligned}}
\MethodTitle{第6小题：$4.2+2.99$}

\MethodBlock{口算思路一：按相同计数单位计算}{先看成若干个百分之一，按整数加减，最后还原成小数。}{\begin{aligned}[t]&\phantom{={}}4.2+2.99\\
&=(420+299)\div 100\\
&=719\div 100\\
&=7.19\end{aligned}}

\MethodBlock{口算思路二：凑整再减去多加的数}{把第二个数凑成整数，再减去多加的部分。}{\begin{aligned}[t]&\phantom{={}}4.2+2.99\\
&=4.2+3-0.01\\
&=7.2-0.01\\
&=7.19\end{aligned}}
\MethodTitle{第7小题：$0.7-0.07$}

\MethodBlock{口算思路一：按相同计数单位计算}{先看成若干个百分之一，按整数加减，最后还原成小数。}{\begin{aligned}[t]&\phantom{={}}0.7-0.07\\
&=(70-7)\div 100\\
&=63\div 100\\
&=0.63\end{aligned}}
\MethodTitle{第8小题：$6.82-4.52$}

\MethodBlock{口算思路一：按相同计数单位计算}{先看成若干个百分之一，按整数加减，最后还原成小数。}{\begin{aligned}[t]&\phantom{={}}6.82-4.52\\
&=(682-452)\div 100\\
&=230\div 100\\
&=2.3\end{aligned}}

\MethodBlock{口算思路二：补差法}{从减数先补到一个整数，再补到被减数，两段差相加。}{\begin{aligned}[t]&\phantom{={}}6.82-4.52\\
&=(5-4.52)+(6.82-5)\\
&=0.48+1.82\\
&=2.3\end{aligned}}
\MethodTitle{第9小题：$0.5+0.51$}

\MethodBlock{口算思路一：按相同计数单位计算}{先看成若干个百分之一，按整数加减，最后还原成小数。}{\begin{aligned}[t]&\phantom{={}}0.5+0.51\\
&=(50+51)\div 100\\
&=101\div 100\\
&=1.01\end{aligned}}
\MethodTitle{第10小题：$2.7+5.2$}

\MethodBlock{口算思路一：按相同计数单位计算}{先看成若干个十分之一，按整数加减，最后还原成小数。}{\begin{aligned}[t]&\phantom{={}}2.7+5.2\\
&=(27+52)\div 10\\
&=79\div 10\\
&=7.9\end{aligned}}
\MethodTitle{第11小题：$0.43-0.24$}

\MethodBlock{口算思路一：按相同计数单位计算}{先看成若干个百分之一，按整数加减，最后还原成小数。}{\begin{aligned}[t]&\phantom{={}}0.43-0.24\\
&=(43-24)\div 100\\
&=19\div 100\\
&=0.19\end{aligned}}
\MethodTitle{第12小题：$4.02-2.7$}

\MethodBlock{口算思路一：按相同计数单位计算}{先看成若干个百分之一，按整数加减，最后还原成小数。}{\begin{aligned}[t]&\phantom{={}}4.02-2.7\\
&=(402-270)\div 100\\
&=132\div 100\\
&=1.32\end{aligned}}

\MethodBlock{口算思路二：补差法}{从减数先补到一个整数，再补到被减数，两段差相加。}{\begin{aligned}[t]&\phantom{={}}4.02-2.7\\
&=(3-2.7)+(4.02-3)\\
&=0.3+1.02\\
&=1.32\end{aligned}}


\par\noindent

\QHeading{第2题　求近似数。}

\PlainItem{$3.951\approx${\color{answerred}$4.0$}（精确到十分位。）}
\par

\PlainItem{$8.839\approx${\color{answerred}$8.84$}（保留两位小数。）}
\PlainItem{$2.9973\approx${\color{answerred}$3.00$}（精确到百分位。）}
\PlainItem{$978694320\approx${\color{answerred}$9.79$}亿（保留两位小数。）}
\par\noindent

\QHeading{第3题　在（　　）里填上适当的数。}

\PlainItem{$6$元$8$分$=$$(${\color{answerred}$6.08$}$)$元}
\par

\Equation{8400$克$=}
\PlainItem{$(${\color{answerred}$8.4$}$)$千克}
\PlainItem{$3.9$米$=$$(${\color{answerred}$390$}$)$厘米}
\Equation{9.04$千克$=}
\PlainItem{$(${\color{answerred}$9$}$)$千克$(${\color{answerred}$40$}$)$克}
\PlainItem{$5.6$千米$=$$(${\color{answerred}$5600$}$)$米}
\Equation{5.4$平方分米$=}
\PlainItem{$(${\color{answerred}$540$}$)$平方厘米}
\PlainItem{$4$吨$70$千克$=$$(${\color{answerred}$4.07$}$)$吨}
\Equation{5$平方米$80$平方分米$=}
\PlainItem{$(${\color{answerred}$5.8$}$)$平方米}
\par\noindent

\Answer{
\QHeading{计算过程：3. 单位换算}
\begin{enumerate}
\item 6元8分换算成元：
\Formula{\begin{aligned}[t]&\phantom{={}}6+8\div100\\&=6+0.08\\&=6.08\text{（元）}\end{aligned}}

\item 3.9米换算成厘米：
\Formula{\begin{aligned}[t]&\phantom{={}}3.9\times100\\&=390\text{（厘米）}\end{aligned}}

\item 5.6千米换算成米：
\Formula{\begin{aligned}[t]&\phantom{={}}5.6\times1000\\&=5600\text{（米）}\end{aligned}}

\item 4吨70千克换算成吨：
\Formula{\begin{aligned}[t]&\phantom{={}}4+70\div1000\\&=4+0.07\\&=4.07\text{（吨）}\end{aligned}}

\item 8400克换算成千克：
\Formula{\begin{aligned}[t]&\phantom{={}}8400\div1000\\&=8.4\text{（千克）}\end{aligned}}

\item 9.04千克中的0.04千克换算成克：
\Formula{\begin{aligned}[t]&\phantom{={}}0.04\times1000\\&=40\text{（克）}\end{aligned}}
所以9.04千克是9千克40克。
\item 5.4平方分米换算成平方厘米：
\Formula{\begin{aligned}[t]&\phantom{={}}5.4\times100\\&=540\text{（平方厘米）}\end{aligned}}

\item 5平方米80平方分米换算成平方米：
\Formula{\begin{aligned}[t]&\phantom{={}}5+80\div100\\&=5+0.8\\&=5.8\text{（平方米）}\end{aligned}}

\end{enumerate}

}

\QHeading{第4题　解决问题。}

\Problem{① 滨海路小学低年级有375名学生，中年级有360名学生，高年级有325名学生。滨海路小学全校共有多少名学生？}
\par


\MethodTitle{方法一}
\Answer{{\Formula{\begin{aligned}[t]&\phantom{={}}375+360+325\\&=375+325+360\\&=(375+325)+360\\&=700+360\\&=1060\text{（名）}\end{aligned}}}\par 答：全校共有1060名学生。}

\MethodTitle{方法二：按先后顺序逐次相加}
\Answer{先合计前两部分，再加第三部分；这是与凑整法对照的常规算法。}
\Answer{\Formula{\begin{aligned}[t]&\phantom{={}}375+360+325\\
&=735+325\\
&=1060\text{（名）}\end{aligned}}}
\Answer{答：全校共有1060名学生。}
\Problem{② 小刚的身高是1.45米，小明比小刚高0.15米，小冬比小明矮0.05米。小冬的身高是多少米？}

\MethodTitle{方法一：简便计算}
\Answer{{\Formula{\begin{aligned}[t]&\phantom{={}}1.45+0.15-0.05\\&=1.45+(0.15-0.05)\\&=1.45+0.1\\&=1.55\text{（米）}\end{aligned}}}\par 答：小冬身高1.55米。}

\MethodTitle{方法二：先求两次高矮变化的净结果}
\Answer{先高0.15米，再矮0.05米，合起来高0.1米。}
\Answer{\Formula{\begin{aligned}[t]&\phantom{={}}1.45+(0.15-0.05)\\
&=1.45+0.1\\
&=1.55\text{（米）}\end{aligned}}}
\Answer{答：小冬身高1.55米。}

